\chapter*{前言}
\addcontentsline{toc}{chapter}{前言}

许多量化研究者第一次需要 \cpp，通常源于系统边界的变化，并非 \python{} “不够好”：研究原型开始处理更大的数据，回测需要可重复的延迟，策略要接入已有的交易基础设施，或者团队希望让资源使用和失败方式变得可预测。本书关注的正是从“能用 \python{} 表达策略”到“能用 \cpp{} 交付可靠量化组件”之间的工程距离。

本书假设你已经熟悉变量、函数、类、容器、异常、测试和虚拟环境。我们不会把一章花在解释循环是什么，而会追问：对象究竟在哪里，谁负责释放资源，复制发生在何时，接口承诺了什么，以及性能结论如何被测量证明。

贯穿全书的项目是一台小型事件驱动回测引擎。它不追求框架功能数量，而追求边界清楚：行情成为事件，策略产生订单意图，撮合器产生成交，组合模块更新状态，统计模块报告结果。每增加一个语言特性，都必须让这个系统更正确、更容易理解或更容易测量。

建议边读边运行代码。遇到性能章节时，先记录假设，再运行基准；遇到并发章节时，先画出所有权和停止协议，再写线程。Modern C++ 的价值不在特性清单，而在于用一套工具让成本、生命周期和契约变得可见。
